\section{Discussion}
\label{sec:discussion}

\paragraph{Position against prior art.}
The heterogeneous budget of Sections~\ref{sec:hetero}--\ref{sec:corr} takes the
scalar timing PL of the companion \citep{baweja2026cti}---itself the rank-one
specialization of MHSS solution separation \citep{blanch2015,joerger2014}---and
makes its per-source bias a constructed, coupled quantity rather than a given,
independent one. Against Gioia's GNSS-only residual T-RAIM \citep{gioia2023} the
step is heterogeneity and traceability; against the all-source solution separation
of Due{\~n}as Arana and colleagues \citep{duenasarana2020} the step is that our
estimand is a scalar time offset with a $\utck{}$ traceability budget, not a
position state; against IEEE~802.1ASed \citep{ieee8021ased} the step is that we emit
an integrity-risk-budgeted protection quantity and, in Section~\ref{sec:gls}, a
monitor for the exact fault selection cannot see---a shift shared by every
candidate. We claim, to our knowledge, that no prior work composes a
source-agnostic bias budget with a common-mode layer as one object; we do not claim
a new estimator, and the whitening \citep{aitken1935} and the uncertainty framework
\citep{jcgm100} are cited as the mature tools they are.

\paragraph{Limitations, stated plainly.}
Four limitations bound this work, and three of them are the recurring caveats.
\emph{(I3)} The correlated-bias theorem is a warning against a variance-style bias
allocation and a motivation for the companion's linear treatment; it does not
harden the companion protection level, which already fuses bias linearly and thus
at least as conservatively (Remark~\ref{rem:i3}). \emph{(I2)} The common-mode
statistic requires an independent external anchor; on the post-fit $H=\bmone_N$
residual it is identically zero, and without an anchor a common shift is
unobservable (Remark~\ref{rem:i2}). \emph{(I1)} The irreducible common-mode blind
spot is large but finite, and our ceiling is a conservative lower bound, not a
total blind spot (Remark~\ref{rem:i1}). Beyond these: the deep integrity tail
of the bias overbound~\eqref{eq:overbound} is \Modelled{}, an extrapolation of a
$95\%$ metrological input that a monthly bulletin cannot validate
(Section~\ref{sec:tracescope}); and the single-realizer-group scenario of
Section~\ref{sec:demo} is a \Modelled{} illustration, not a survey of real shared
realizers. No result here is externally \Validated{}; the only \Validated{} row in
the arc is the companion's holdover overbound coverage, and the BIPM Circular-T
series that sizes our biases is a \Cited{} input, not re-validated
(Section~\ref{sec:datasetnote}).

\paragraph{What the rest of the arc adds.}
This paper is P2 of a four-paper Composed Timing Integrity arc and is honest about
what it defers. The adversary---delay attacks, which are provably undetectable in
the symmetric single-path case \citep{mizrahi7384,narula2018} and are attacked
protocol-specifically by monitors such as PTPsec \citep{finkenzeller2024ptpsec}---is
the subject of the adversarial companion P3; here the fault model is benign
common-mode and shared-realizer bias, and the anchor of caveat~I2 is assumed
trustworthy. The question of whether any monitor, including this one, is
\emph{sound} on a common rule---mapping its coverage rather than flattering
it---is the subject of the open benchmark companion P4, which scores methods on a
shared fault menu and makes no accuracy claim by design. This paper contributes the
one missing composed object, and its honest blind-spot characterization, between
them.
